\chapter{Thích Ngồi Cuối Lớp}
\chapterintrobi{Tầm nhìn rộng nhưng mức tập trung thường hẹp}{Love the back row}{Ngồi cuối lớp không có tội, nhưng nhiều bạn xem đó như vùng khí hậu riêng để giảm áp lực học hành.}

\begin{lucbatbox}{Lục bát: Hậu phương bình yên}
Cuối lớp gió mát hơn nhiều\
Xa tầm gọi hỏi, lại chiều lòng em\
Từ đây quan sát rất êm\
Bạn nào bị gọi đứng lên thấy liền\
Chỉ là kiến thức hơi nghiêng\
Đi từ bảng xuống tới riêng chỗ ngồi
\end{lucbatbox}

\begin{tutuyetbox}{Thất ngôn tứ tuyệt: Vùng an toàn}
Ngồi cuối lớp nghe có vẻ hay\
Ít khi thầy gọi tới nơi này\
Nhưng điều xa nhất không là bục giảng\
Mà là chú ý cứ vơi đầy
\end{tutuyetbox}

\begin{ghichu}
Ngồi cuối lớp không làm điểm tự rớt. Nhưng nếu đi cùng thói quen “ẩn mình”, nó thường kéo theo nhiều thứ khác rơi nhanh hơn.
\end{ghichu}

\begin{englishnote}
\textbf{back row}: hàng cuối.\par
\textbf{safe zone}: vùng an toàn.\par
\textbf{I like sitting at the back.}: Em thích ngồi cuối lớp.
\end{englishnote}